\documentclass[11pt]{article}
\usepackage[utf8]{inputenc}
\usepackage[T1]{fontenc}
\usepackage[spanish]{babel}
\usepackage[margin=2.3cm]{geometry}
\usepackage{amsmath,amssymb}
\usepackage{enumitem}
\setlength{\parindent}{0pt}
\setlength{\parskip}{6pt}

\begin{document}

\begin{center}
{\Large\bfseries Control 1 --- Microeconomía II}\\[4pt]
{\large Simulacro (formato Mundaca)}\\[4pt]
80 puntos \quad$\cdot$\quad 80 minutos \quad$\cdot$\quad sin apuntes
\end{center}

\hrule
\medskip

\textbf{I. Comentes (5 pts c/u).} Comente si las siguientes afirmaciones son verdaderas, falsas o inciertas. Justifique su respuesta; la nota depende de la justificación.

\begin{enumerate}[leftmargin=*, itemsep=8pt]
\item En la discriminación de precios de 2.\textsuperscript{o} grado, el monopolista ofrece al consumidor de baja demanda la misma cantidad que le ofrecería en competencia perfecta, al igual que al consumidor de alta demanda.

\item Con una tarifa en dos partes y consumidores idénticos, el monopolista fija el precio por unidad por sobre el costo marginal, porque así extrae más excedente del consumidor.

\item En el equilibrio de largo plazo de competencia monopolística, cada empresa produce en el mínimo de su costo medio, al igual que en competencia perfecta.

\item Un monopolista que discrimina precios de 3.\textsuperscript{er} grado cobra un precio más alto en el mercado cuya demanda es más elástica.
\end{enumerate}

\textbf{II. Desarrollo (10 pts c/u).}

\begin{enumerate}[leftmargin=*, itemsep=8pt, resume]
\item ¿Qué es un monopolio natural? Compare la regulación que fija el precio igual al costo marginal con la que lo fija igual al costo medio, en términos de eficiencia y de financiamiento de la empresa. ¿Qué problema práctico enfrenta el regulador al aplicar cualquiera de las dos? Grafique.

\item Un monopolista no puede distinguir a los consumidores de alta y de baja demanda. Explique cómo puede, de todos modos, discriminar precios, y por qué le conviene distorsionar la cantidad que ofrece al consumidor de baja demanda.
\end{enumerate}

\textbf{III. Matemático (40 pts).}

Un monopolista posee una función de costos totales de la forma $c(q) = q^{2} + 20q + 200$ y abastece a un mercado cuya función de demanda es $q(p) = 100 - \dfrac{p}{2}$.

\begin{enumerate}[label=\alph*), leftmargin=*, itemsep=8pt]
\item \textbf{(20 pts)} Obtenga todos los elementos del equilibrio del monopolista.

\item \textbf{(20 pts)} Calcule la pérdida de eficiencia que sufre esta economía. En específico, señale cómo varían el precio y la cantidad de equilibrio y los excedentes del consumidor, del productor y social respecto de la situación de competencia perfecta. ¿Produce el monopolista la mitad de la cantidad de competencia perfecta? Explique. Grafique ambos apartados.
\end{enumerate}

\medskip
\hrule
\smallskip
{\small\textit{Antes de empezar, los 4 chequeos:} (1) reparte el tiempo por puntaje (matemático 40 min, desarrollos 10 c/u, comentes 5 c/u); (2) subraya el sujeto de cada pregunta; (3) cierra cada «compare» con la comparación explícita; (4) no reuses cifras de otro ítem sin revisarlas. Veredicto V/F/I en la primera palabra; trabaja con fracciones.}

\end{document}
